% Standalone problem document, generated from the adjacent README.md.
% Compile with XeLaTeX. Mathematical content is maintained in Markdown.
\documentclass[11pt,a4paper]{article}
\usepackage[fontsize=10.5pt]{scrextend}
\usepackage[margin=22mm,headheight=15pt,headsep=9mm,footskip=11mm]{geometry}
\usepackage{fontspec}
\setmainfont{texgyrepagella}[Extension=.otf,UprightFont=*-regular,BoldFont=*-bold,ItalicFont=*-italic,BoldItalicFont=*-bolditalic]
\setsansfont{texgyreheros}[Extension=.otf,UprightFont=*-regular,BoldFont=*-bold,ItalicFont=*-italic,BoldItalicFont=*-bolditalic]
\setmonofont{texgyrecursor}[Extension=.otf,UprightFont=*-regular,BoldFont=*-bold,ItalicFont=*-italic,BoldItalicFont=*-bolditalic,Scale=MatchLowercase]
\usepackage{amsmath,amssymb}
\usepackage{unicode-math}
\setmathfont{texgyrepagella-math.otf}
\usepackage{xcolor,microtype,enumitem,needspace,fancyhdr}
\usepackage{bookmark}
\definecolor{ink}{HTML}{17344B}
\definecolor{accent}{HTML}{197D80}
\definecolor{muted}{HTML}{596773}
\hypersetup{colorlinks=true,linkcolor=accent,urlcolor=accent,pdftitle={IS-02 - Where
a symmetric stochastic matrix can be spectrally
unique},pdfauthor={Open Problems in Numerical Linear Algebra}}
\urlstyle{same}
\setlength{\parindent}{0pt}
\setlength{\parskip}{5pt plus 1pt minus 1pt}
\linespread{1.04}
\setlength{\emergencystretch}{3em}
\widowpenalty=10000
\clubpenalty=10000
\displaywidowpenalty=10000
\allowdisplaybreaks[1]
\setlist{nosep,leftmargin=1.5em}
\providecommand{\tightlist}{\setlength{\itemsep}{0pt}\setlength{\parskip}{0pt}}
\providecommand{\pandocbounded}[1]{#1}
\pagestyle{fancy}
\fancyhf{}
\fancyhead[L]{\sffamily\footnotesize\color{muted}OPEN PROBLEMS IN NUMERICAL LINEAR ALGEBRA}
\fancyhead[R]{\sffamily\bfseries\footnotesize\color{accent}IS-02}
\fancyfoot[L]{\parbox[b]{0.9\textwidth}{\sffamily\fontsize{8}{10}\selectfont\color{muted}Editorial ratings; a bounded search cannot certify that a problem remains unsolved.\\Literature check: 2026-09-10}}
\fancyfoot[R]{\sffamily\footnotesize\color{muted}\thepage}
\renewcommand{\headrulewidth}{0.3pt}
\renewcommand{\headrule}{\hbox to\headwidth{\color{accent}\leaders\hrule height \headrulewidth\hfill}}
\makeatletter
\renewcommand{\section}{\Needspace{5\baselineskip}\@startsection{section}{1}{0pt}{12pt plus 2pt minus 2pt}{4pt}{\sffamily\bfseries\large\color{ink}}}
\renewcommand{\subsection}{\Needspace{4\baselineskip}\@startsection{subsection}{2}{0pt}{9pt plus 1pt minus 1pt}{3pt}{\sffamily\bfseries\normalsize\color{ink}}}
\makeatother
\setcounter{secnumdepth}{0}
\begin{document}
{\sffamily\small\color{muted}Eigenvalues and inverse problems\par}
\vspace{3pt}
{\raggedright\hyphenpenalty=10000\exhyphenpenalty=10000\sffamily\bfseries\fontsize{21}{24}\selectfont\color{ink}Where
a symmetric stochastic matrix can be spectrally unique\par}
\vspace{7pt}
{\sffamily\small\textbf{Difficulty:} challenging\quad\textbf{Importance:} interesting
to specialist\par}
{\sffamily\small\bfseries\color{accent}Status: Open\par}
\vspace{4pt}
{\color{accent}\hrule height 0.8pt}
\vspace{6pt}
\textbf{Rating rationale:} Challenging reflects the geometry of
isospectral stochastic families in arbitrary dimension; specialist
impact concerns the narrow property of spectral uniqueness within that
class.

\subsection{Problem statement}\label{problem-statement}

For \(n\geq4\), put
\(\mathcal S_n=\{A\in\mathbb R^{n\times n}:A=A^T,\ A\geq0,\ A\mathbf1=\mathbf1\}\)
and \(C_n=(\mathbf1\mathbf1^T-I)/(n-1)\). A matrix \(A\in\mathcal S_n\)
is \emph{spectrally unique} if every \(B\in\mathcal S_n\) with the same
eigenvalues, including multiplicities, satisfies \(B=R^TAR\) for a
permutation matrix \(R\). Let \([X,Y]=\{(1-t)X+tY:0\leq t\leq1\}\).

Prove or disprove the following necessary condition: every spectrally
unique \(A\in\mathcal S_n\) with \(\operatorname{tr}A>0\) belongs to

\[
[I,C_n]\ \cup\!
\bigcup_{V\in\operatorname{vert}(\mathcal S_n)}
\bigl([I,V]\cup[C_n,V]\bigr).
\]

Here a vertex is an extreme point of the indicated convex polytope; it
need not be a permutation matrix. Only the stated implication is
asserted. This asks where recovering a nonnegative symmetric stochastic
matrix from its spectrum can be unique up to relabeling, a different
issue from mere spectral feasibility.

\subsection{References}\label{references}

Mourad and Abbas, \href{https://arxiv.org/pdf/1310.1273}{2013 preprint},
definitions in §1 and Conjecture 5.1, p.~10;
\href{https://doi.org/10.1080/03081087.2014.903590}{published article},
\emph{Linear and Multilinear Algebra} 63 (2015), 869--881.

\subsection{Earlier status check ---
2026-09-08}\label{earlier-status-check-2026-09-08}

Searches for the exact title,
\textit{symmetric doubly stochastic Conjecture 5.1}, and
author/title combinations with \textit{counterexample} and \textit{2026}
found no resolution. The source solves order three, which is excluded
from the remaining statement above.

\subsection{Audit update --- 2026-09-10}\label{audit-update-2026-09-10}

Rechecked \href{https://arxiv.org/pdf/1310.1273}{Mourad--Abbas}, §5,
Conjecture 5.1, against the trace and locus restrictions here. Searches
for spectrally unique symmetric stochastic matrices and later work on
that conjecture found no general answer. The order-three classification
is outside the displayed unresolved dimensions, so it does not change
this entry's status. Evidence remains historical rather than a recent
explicit reaffirmation.
\end{document}
